\begin{env}[6.15.3]
\label{IV.2.s6.15.3-fr}
Diremos, abreviadamente, que un morfismo $f:X\to Y$ es
\emph{radicial en un punto} $y\in Y$ si $f^{-1}(y)$ es vacío
o consta de un único punto $x$ y si $\boldsymbol{k}(x)$ es una extensión
radicial de $\boldsymbol{k}(y)$. Por tanto, decir que $f$ es radicial en
todos los puntos de $Y$ equivale a decir que $f$ es radicial
(\textbf I, 3.5.8). Si $g:f^{-1}(y)\to\operatorname{Spec}(\boldsymbol{k}(y))$ es el
morfismo obtenido de $f$ por el cambio de base
$\operatorname{Spec}(\boldsymbol{k}(y))\to Y$, entonces $f$ es radicial en el
punto $y\in Y$ si y solo si $g$ es \emph{radicial}.
\end{env}

\begin{lemma}[6.15.3.1]
\label{IV.2.s6.15.3.1-fr}
\begin{enumerate}
\item[\textup{(i)}] Sean $f:X\to Y$ y $g:Y\to Z$ dos morfismos. Si $f$
es radicial en un punto $y$ y $g$ es radicial en el punto
$z=g(y)$, entonces $g\circ f$ es radicial en el punto $g(y)$.
La recíproca es cierta si $f$ es sobreyectivo.
\item[\textup{(ii)}] Sean $f:X\to Y$ y $h:Y'\to Y$ dos morfismos, y sea
$f'=f_{(Y')}:X\times_Y Y'\to Y'$. Sean $y'\in Y'$ e $y=h(y')$;
entonces $f$ es radicial en el punto $y$ si y solo si $f'$ es
radicial en el punto $y'$.
\end{enumerate}
\end{lemma}
